\begin{thebibliography}{99}

\bibitem{Sernau2026}
L. Sernau and collaborators,
\emph{GRH for Dirichlet $L$-functions, in Lean 4},
repository \code{lsernau/grh-lean}, commit
\code{72fead102f663d5c7a407964a13dc0cf9d977207}, 25 November 2026.
\url{https://github.com/lsernau/grh-lean/tree/72fead102f663d5c7a407964a13dc0cf9d977207}.

\bibitem{ProveIt2026}
V. Reshetnikov, \emph{Generalized Riemann Hypothesis: Article},
ProveIt repository, commit
\code{22e5c8b7dda558fd9a50b61465483cb4c0433694}, inspected 27 November 2026.
In particular \code{sections/sel.tex} and \code{sections/ana.tex}.
\url{https://github.com/VladimirReshetnikov/ProveIt/tree/22e5c8b7dda558fd9a50b61465483cb4c0433694/Analysis/GeneralizedRiemannHypothesis/Article}.

\bibitem{Sutherland2026}
A. V. Sutherland,
\emph{Verification files for the shorter Lean proof of RH and of GRH
for Dirichlet $L$-functions}, updated 27 November 2026.
\url{https://math.mit.edu/~drew/shortgrh/}.

\bibitem{Carayol1994}
H. Carayol,
\emph{Formes modulaires et repr\'esentations galoisiennes \`a valeurs
dans un anneau local complet},
Contemporary Mathematics \textbf{165} (1994), 213--237.
\url{https://doi.org/10.1090/conm/165/01601}.

\bibitem{Serre1993}
J.-P. Serre, notes by R. Rouquier,
\emph{Repr\'esentations lin\'eaires sur des anneaux locaux, d'apr\`es
Carayol}, lecture of 8 October 1993, Cours de Jean-Pierre Serre,
volume 14. Theorem 1 is the classical trace-rigidity statement.
\url{https://www.numdam.org/item/CJPS_1993__14_.pdf}.

\bibitem{Brumer2019}
A. Brumer, A. Pacetti, C. Poor, G. Tornar\'ia, J. Voight and D. S. Yuen,
\emph{On the paramodularity of typical abelian surfaces},
Algebra \& Number Theory \textbf{13} (2019), 1145--1195.
\url{https://doi.org/10.2140/ant.2019.13.1145}.
Used as related context for trace-based representation comparison,
not as an input to the new bound.

\bibitem{Hogenauer1981}
E. B. Hogenauer,
\emph{An Economical Class of Digital Filters for Decimation and
Interpolation}, IEEE Transactions on Acoustics, Speech, and Signal
Processing \textbf{29} (1981), no. 2, 155--162.
\url{https://doi.org/10.1109/TASSP.1981.1163535}.
Equation (3) gives the classical cascade polynomial used here.

\bibitem{Bennett2021}
M. A. Bennett, G. Martin, K. O'Bryant and A. Rechnitzer,
\emph{Counting zeros of Dirichlet $L$-functions},
Mathematics of Computation \textbf{90} (2021), 1455--1482.
\url{https://doi.org/10.1090/mcom/3599}.
Author copy:
\url{https://personal.math.ubc.ca/~bennett/BeMaObRe-MathComp-2021.pdf}.

\end{thebibliography}
